\bookchapter{Kapitel 65: Projekt 1 – Mein Haushaltsbuch (Lokalversion)}

\chapterintro{Willkommen bei deinem ersten „richtigen Projekt“. Dieses Haushaltsbuch kommt zwar ohne Netz aus, ist aber komplett ausgestattet: Buchungen erfassen, Verlauf ansehen, monatliche Auswertung, Löschen und Ändern. Wenn du fertig bist, hast du den kompletten Ablauf „Daten Logik Oberfläche“ einmal durchlaufen. Wilde nicht gleich den Code ab – erst durchdenken, dann anpacken; das ist unser bewährtes Prinzip.}

\booksection{Anforderungen überdenken: erst fünf Fragen stellen}

\begin{center}
\begin{tabularx}{\linewidth}{|X|X|}
\hline
\textbf{Element} & \textbf{Antwort} \tabularnewline
\hline
Who (für wen?) & Für dich selbst; Familie und Freunde gehen auch \tabularnewline
What (was?) & Einnahmen/Ausgaben erfassen, Verlauf ansehen, Monatsbilanz, Auswertung nach Kategorien \tabularnewline
Why (warum?) & Übung + herausfinden, wohin das Geld geht \tabularnewline
How (wie?) & Swift + SwiftUI + SwiftData + MVVM, rein lokal ohne Backend \tabularnewline
Constraint (Randbedingungen) & In einer Woche fertig; mindestens iOS 17 (SwiftData-Voraussetzung) \tabularnewline
\hline
\end{tabularx}
\end{center}

\begin{center}\small\color{gray}Die fünf Anforderungselemente des Haushaltsbuchs\end{center}

Dann zerlegst du das Ganze in kleine Schritte, nach denen jeweils etwas Benutzbares entsteht: 1) die SwiftData-Datenschicht bauen (eine Buchung speichern können); 2) die Listenseite (Verlauf ansehen); 3) den Eingabe-Dialog (Buchung hinzufügen); 4) Löschen und Bearbeiten; 5) die Monatsbilanz-Karte. Nach jedem erledigten Schritt ist die App ein „brauchbares Ding“.

\booksection{Die Datenschicht: den „Buchungen“ ein Zuhause geben}

Definiere zuerst, wie eine Buchung aussieht – die Felder sollen ausreichen, aber nicht gierig sein:

\begin{lstlisting}[language=Swift]
import SwiftData

@Model
class Record {
    var type: String      // "income" Einnahme | "expense" Ausgabe
    var amount: Double    // Betrag (Yuan)
    var category: String  // Kategorie: Essen, Verkehr, Gehalt …
    var note: String      // Notiz (optional)
    var date: Date        // Buchungsdatum

    init(type: String, amount: Double, category: String, note: String = "", date: Date = Date()) {
        self.type = type
        self.amount = amount
        self.category = category
        self.note = note
        self.date = date
    }
}
\end{lstlisting}

\begin{tcbinfo}{Tipp: Warum Double für den Betrag?}
Seröse Finanz-Apps verwenden Decimal (exakte Dezimalzahlen), um Rundungsfehler bei Gleitkommazahlen zu vermeiden (das alte Problem, dass 0,1 + 0,2 $\neq$ 0,3). Für ein Übungsprojekt reicht Double völlig; baust du eines Tages ein echtes Finanzprodukt, denk an den Wechsel zu Decimal oder ein „Int in Cent“ – eine Falle, in die unzählige Vorgänger getappt sind.

\end{tcbinfo}

\booksection{MVVM: Oberfläche und Logik „ziehen getrennte Wege“}

Die Oberfläche direkt auf die Datenbank zugreifen lassen? Läuft zwar, aber der Oberflächencode wird lang und unappetitlich. Die MVVM-Architektur sorgt für klare Zuständigkeit: View (Oberfläche) zeichnet nur, ViewModel (Logikzentrale) kümmert sich um die Daten, Model (SwiftData) um die Speicherung.

View (SwiftUI-Oberfläche)

ViewModel (@Observable als Logik-Zentrale)

Model (Datenlayer mit SwiftData)

\begin{center}\small\color{gray}Abbildung 4: Die drei Ebenen von MVVM – Daten fließen nur in eine Richtung, und keine Ebene greift einer anderen ins Handwerk\end{center}

\begin{lstlisting}[language=Swift]
import Observation

@Observable
class RecordViewModel {
    var records: [Record] = []
    private let context: ModelContext

    init(context: ModelContext) { self.context = context }

    func load() {
        // Filtern mit #Predicate von SwiftData + Sortierung
        let descriptor = FetchDescriptor<Record>(
            sortBy: [SortDescriptor(\.date, order: .reverse)]
        )
        records = (try? context.fetch(descriptor)) ?? []
    }

    func add(type: String, amount: Double, category: String, note: String) {
        let r = Record(type: type, amount: amount, category: category, note: note)
        context.insert(r)
        try? context.save()
        load()                       // neu laden; die Oberfläche aktualisiert sich
    }

    func delete(_ record: Record) {
        context.delete(record)
        try? context.save()
        load()
    }
}
\end{lstlisting}

In der Oberfläche @State private var vm = RecordViewModel(...) anlegen und die Liste direkt aus vm.records rendern – ändern sich die Daten, aktualisiert sich die Oberfläche automatisch.

\booksection{Die Oberfläche: Listenseite + Eingabe-Dialog}

\begin{lstlisting}[language=Swift]
struct RecordListView: View {
    @State private var vm: RecordViewModel
    @State private var showAdd = false

    var body: some View {
        NavigationStack {
            List {
                Section("Saldo des Monats") {          // Zusammenfassungs-Karte
                    SummaryRow(records: vm.records)
                }
                Section("Umsätze") {
                    ForEach(vm.records) { r in
                        RecordRow(record: r)
                    }
                    .onDelete { idx in
                        idx.forEach { vm.delete(vm.records[$0]) }
                    }
                }
            }
            .navigationTitle("Mein Ausgabenbuch")
            .toolbar {
                Button { showAdd = true } label: {
                    Image(systemName: "plus")
                }
            }
            .sheet(isPresented: $showAdd) {
                AddRecordView(vm: vm)          // Eingabedialog
            }
        }
    }
}
\end{lstlisting}

Wichtige Punkte für den Eingabe-Dialog: für das Betragsfeld .keyboardType(.decimalPad) (Zifferntastatur) verwenden, für Einnahmen/Ausgaben ein segmentiertes Steuerelement (Picker), die Kategorie über ein Menü wählen, und bei „Speichern“ vm.add(...) aufrufen und den Dialog mit dismiss() schließen.

\begin{tcbwarn}{Achtung: Bei der Eingabeprüfung nicht faulenzen}
Das Betragsfeld musst du unbedingt prüfen: Leere Werte, negative Zahlen und unzulässige Eingaben wie “abc” müssen abgefangen werden (mit Double(text) prüfen; im Fehlerfall kommt nil zurück). Einer der häufigsten Gründe für App-Abstürze ist, „die Eingaben der Nutzer für brav zu halten“. Denk dran: Alles, was Nutzer eingeben, musst du wie einen Gegner unter die Lupe nehmen.

\end{tcbwarn}

\booksection{Bearbeiten, Löschen und Bilanz}

\begin{itemize}
  \item Löschen: .onDelete in List fürs Löschen per Linkswischen + Bestätigungsdialog (versehentliches Löschen ist das, was Nutzer von Buchhaltungs-Apps am meisten hassen);
  \item Bearbeiten: ein Tipp auf die Zeile öffnet den Bearbeitungsdialog mit vorausgefüllten Originalwerten (ein Wort ändern ist bequemer als alles neu eintippen);
  \item Bilanz: mit filter + reduce „Einnahmen / Ausgaben / Saldo des Monats“ berechnen und oben in einer Karte anzeigen:
\end{itemize}

\begin{lstlisting}[language=Swift]
func monthSummary(_ records: [Record]) -> (income: Double, expense: Double) {
    let cal = Calendar.current
    let monthRecords = records.filter {
        cal.isDate($0.date, equalTo: Date(), toGranularity: .month)
    }
    let income = monthRecords
        .filter { $0.type == "income" }
        .reduce(0) { $0 + $1.amount }
    let expense = monthRecords
        .filter { $0.type == "expense" }
        .reduce(0) { $0 + $1.amount }
    return (income, expense)
}
\end{lstlisting}

\booksection{Abnahme: Prüfung auf drei Ebenen + zwei Zusatzaufgaben für Fortgeschrittene}

Nach Fertigstellung prüfst du das Projekt auf drei Ebenen entlang von „Funktion Grenzfälle Architektur“: Lässt sich erfassen, lässt sich löschen, sind die Daten nach einem Neustart noch da (Funktion)? Stürzt die App ab, wenn du 0 oder negative Beträge eingibst (Grenzfälle)? Sind Oberfläche und Daten vermischt, sind Felder fest einprogrammiert (Architektur)?

Nach bestandener Abnahme der Basisversion gönnst du dir zwei „Fortgeschrittenen-Gerichte“:

\begin{enumerate}
  \item Kategorien-Auswertung: Ergänze einen „Statistik“-Reiter und zeichne mit Swift Charts die Ausgaben je Kategorie als Kreis-/Balkendiagramm – Apples offizielle Diagrammbibliothek, mit wenigen Codezeilen zum fertigen Chart;
  \item iCloud-Sync: Für iCloud-Synchronisation mit SwiftData genügen ein paar Einstellungen – das Haushaltsbuch synct automatisch zwischen iPhone und iPad, und auf einen Schlag wirkt alles „profihaft“.
\end{enumerate}

\begin{tcbpractice}{Jetzt bist du dran: diese Woche fertigstellen}
Arbeite in der Reihenfolge „Datenschicht ViewModel Listenseite Eingabe-Dialog Bilanz-Karte“, jeden Tag ein Schritt. Wenn du hängst, wirf die Fehlermeldung wortwörtlich einer Suchmaschine oder einer KI hin (erinnerst du dich an Teil 7, „KI-Programmierung in der Praxis“? Beschreibe beim Schildern des Problems klar „erwartetes Verhalten, tatsächliches Verhalten, Schritte zur Reproduktion“). Wenn du fertig bist, führe damit eine Woche lang dein Haushaltsbuch – eine selbst gebaute App zu benutzen gibt ein Erfolgserlebnis, das dich weiter antreibt.

\end{tcbpractice}

\begin{tcbsummary}{Kapitelfazit}
\begin{itemize}
  \item Fünf Anforderungselemente in kleine Schritte zerlegen, nach denen jeweils etwas Benutzbares entsteht, dann loslegen
  \item Datenschicht: @Model-Modell + modelContainer-Container + FetchDescriptor-Abfragen
  \item MVVM-Arbeitsteilung: View zeichnet nur, das @Observable-ViewModel kümmert sich um die Logik, SwiftData um die Speicherung
  \item Oberfläche: List + Section + toolbar + sheet als Eingabe-Dialog
  \item Abnahme auf drei Ebenen (Funktion/Grenzfälle/Architektur) + Zusatzaufgaben (Swift Charts, iCloud-Sync)
\end{itemize}

\end{tcbsummary}
